\subsection{Постановка задачи}

Второй микропроект посвящён моделированию распространения короткого
ультразвукового импульса в трёхмерной анатомической модели головы.
Искомая величина~--- акустическое давление $p(\mathbf{x}, t)$,
распространяющееся по неоднородной среде: мягким тканям, костям,
кровеносным сосудам и другим областям, заданным на тетраэдральной сетке.

Цель работы~--- собрать полный расчётный пайплайн:
\begin{enumerate}
  \item взять готовую сегментированную 3D-модель головы;
  \item перенести материал каждой ячейки в DOLFINx;
  \item задать импульсный ультразвуковой источник;
  \item решить нестационарную акустическую задачу методом конечных элементов;
  \item сохранить результаты в формате \texttt{.vtu} для ParaView.
\end{enumerate}

Геометрическая модель взята из репозитория
\href{https://github.com/avasyukov/gcm-3d/tree/markered-corpse/models/ani3d}{gcm-3d},
ветка с моделями \texttt{ani3d}. Это готовая тетраэдральная сетка
анатомической области с метками материалов. Расчётный код был написан
с опорой на готовые примеры из
\href{https://github.com/adeebkor/fenicsx-fus}{fenicsx-fus}, где уже
используется DOLFINx для задач ультразвука.

\subsection{Расчётная модель}

\subsubsection{Геометрия и материалы}

В каталоге \texttt{mini\_project/3d-models/} используются два основных
VTU-файла:
\begin{itemize}
  \item \texttt{mesh.vtu}~--- исходная неструктурированная тетраэдральная
        сетка головы;
  \item \texttt{head\_with\_materials.vtu}~--- сетка после предобработки,
        где к каждой ячейке добавлены поля \texttt{material\_id},
        \texttt{rho}, \texttt{lambda}, \texttt{mu},
        \texttt{sound\_speed}.
\end{itemize}
В итоговой сетке:
\[
  N_{\rm points}=125\,215, \qquad N_{\rm cells}=729\,468.
\]

Координаты модели используются в миллиметрах, время~--- в секундах,
поэтому скорость звука в коде задаётся в мм/с. Плотность и упругие
константы берутся из таблицы материалов, приведённой в том же источнике,
что и модель.

\begin{center}
\begin{tabular}{lcccc}
  \hline
  Область & Material ID & $\rho$, г/мм$^3$ & $c$, мм/с
          & $(\lambda, \mu)$ \\
  \hline
  Жир & 3, 31 & $1{,}158\cdot 10^{-3}$ & $1{,}45\cdot 10^6$
      & $(3{,}07\cdot 10^6,\;2{,}05\cdot 10^6)$ \\
  Мышцы / мягкие ткани & 4, 11, 12, 27, 30, 32
      & $1{,}00\cdot 10^{-3}$ & $1{,}58\cdot 10^6$
      & $(3{,}07\cdot 10^6,\;2{,}05\cdot 10^6)$ \\
  Кость & 21 & $1{,}00\cdot 10^{-3}$ & $3{,}50\cdot 10^6$
      & $(7{,}86\cdot 10^8,\;1{,}18\cdot 10^9)$ \\
  Кровь & 22, 28, 29 & $1{,}00\cdot 10^{-3}$ & $1{,}57\cdot 10^6$
      & $(2{,}39\cdot 10^5,\;5{,}00\cdot 10^2)$ \\
  Трахея & 24 & $2{,}00\cdot 10^{-3}$ & $1{,}54\cdot 10^6$
      & $(1{,}43\cdot 10^7,\;3{,}57\cdot 10^6)$ \\
  Аорта & 25 & $1{,}00\cdot 10^{-3}$ & $1{,}57\cdot 10^6$
      & $(9{,}21\cdot 10^6,\;1{,}90\cdot 10^5)$ \\
  Вены & 26 & $1{,}00\cdot 10^{-3}$ & $1{,}57\cdot 10^6$
      & $(3{,}289\cdot 10^7,\;6{,}70\cdot 10^5)$ \\
  \hline
\end{tabular}
\end{center}

В акустической постановке непосредственно используются $\rho(\mathbf{x})$
и $c(\mathbf{x})$. Константы Ламе $\lambda(\mathbf{x})$, $\mu(\mathbf{x})$
сохранены в модели для следующего шага~--- линейно-упругого расчёта
смещений.

\subsubsection{Предобработка сетки}

Файл \texttt{preprocess.py} читает VTU через \texttt{meshio}, находит
тетраэдральный блок и восстанавливает \texttt{Material ID} для каждой
ячейки.

Материалы записываются в ячейки:
\begin{equation}
  \rho_h,\; c_h,\; \lambda_h,\; \mu_h \in DG_0(\Omega),
\end{equation}
то есть каждое значение постоянно внутри одного тетраэдра.

\subsection{Физическая постановка}

\subsubsection{Линейная акустика}

Ультразвуковой импульс рассматривается как малая добавка давления
к равновесному состоянию среды. Поэтому используется линейная акустика:
амплитуда возмущения мала по сравнению с модулем объёмного сжатия,
а нелинейные эффекты и поглощение в первом приближении не учитываются.

Пусть $\mathbf{v}(\mathbf{x},t)$~--- акустическая скорость частиц среды,
$p(\mathbf{x},t)$~--- акустическое давление, $\rho(\mathbf{x})$~---
плотность, $c(\mathbf{x})$~--- скорость звука. Для малых возмущений:
\begin{equation}
  \rho(\mathbf{x})\,\frac{\partial \mathbf{v}}{\partial t}
  = -\nabla p,
\end{equation}
\begin{equation}
  \frac{1}{\rho(\mathbf{x})c^2(\mathbf{x})}
  \frac{\partial p}{\partial t}
  + \nabla\cdot\mathbf{v} = 0.
\end{equation}
Первое уравнение~--- линеаризованный закон Ньютона: градиент давления
разгоняет частицы среды. Второе~--- связь сжимаемости и давления:
при сжатии объёма давление растёт; коэффициент
$K=\rho c^2$ играет роль объёмного модуля.

Исключая $\mathbf{v}$, получаем волновое уравнение в неоднородной среде:
\begin{equation}\label{eq:acoustic_strong}
  \frac{1}{\rho(\mathbf{x})c^2(\mathbf{x})}
  \frac{\partial^2 p}{\partial t^2}
  - \nabla\cdot\left(
      \frac{1}{\rho(\mathbf{x})}\nabla p
    \right)
  = s(\mathbf{x},t).
\end{equation}

Если $\rho$ и $c$ постоянны, оно сводится к обычному уравнению
$p_{tt}-c^2\Delta p = \rho c^2 s$. В данной задаче коэффициенты
кусочно-постоянны по тканям, поэтому фронт волны меняет скорость,
отражается и преломляется на границах материалов.

\subsubsection{Источник}

Источник разделён на пространственную и временную части:
\begin{equation}
  s(\mathbf{x},t)=g(\mathbf{x})\,q(t).
\end{equation}

Пространственный профиль~--- гауссово пятно:
\begin{equation}
  g(\mathbf{x}) =
  A \exp\left(
    -\frac{\|\mathbf{x}-\mathbf{x}_s\|^2}{2\sigma^2}
  \right).
\end{equation}
В коде центр $\mathbf{x}_s$ ставится около нижней границы расчётной
области:
\begin{equation}
  x_s=\frac{x_{\min}+x_{\max}}{2}, \qquad
  y_s=y_{\min}+0{,}4(y_{\max}-y_{\min}), \qquad
  z_s=z_{\min},
\end{equation}
а ширина берётся как $\sigma=0{,}1\,D$, где $D$~--- максимальный
габаритный размер модели.

Временная форма~--- короткий ультразвуковой пакет с плавным окном:
\begin{equation}\label{eq:burst}
q(t)=
\begin{cases}
  \sin(2\pi f\tau)\,
  \sin^2\left(\dfrac{\pi\tau}{T}\right),
  & 0 \le \tau \le T,\\[8pt]
  0, & \text{иначе},
\end{cases}
\end{equation}
где
\[
  \tau=t-t_{\rm start}, \qquad
  t_{\rm start}=t_0-\frac{T}{2}, \qquad
  T=\frac{N}{f}.
\]

В расчёте:
\[
  f = 1~\text{МГц}, \qquad N=5, \qquad
  t_0=\frac{5}{f}, \qquad T=\frac{5}{f}.
\]
Множитель $\sin^2(\pi\tau/T)$ плавно поднимает амплитуду от нуля и
возвращает её обратно к нулю. Это уменьшает артефакты, которые появились
бы при резком обрыве синуса.

\subsubsection{Граничные условия}

Явно заданных граничных условий Дирихле в акустическом коде нет.
После интегрирования по частям естественное граничное условие имеет вид
\begin{equation}
  \left.
  \frac{1}{\rho}\nabla p\cdot\mathbf{n}
  \right|_{\partial\Omega}=0.
\end{equation}
Физически это отражающая граница.

\subsection{Численный метод}

\subsubsection{Слабая формулировка}

Умножим уравнение~\eqref{eq:acoustic_strong} на тестовую функцию
$v \in V$ и проинтегрируем по области. С учётом естественного
граничного условия получаем:
\begin{equation}\label{eq:acoustic_weak}
  \int_\Omega
  \frac{1}{\rho c^2}\,
  \frac{\partial^2 p}{\partial t^2}\,v\,d\mathbf{x}
  +
  \int_\Omega
  \frac{1}{\rho}\,\nabla p\cdot\nabla v\,d\mathbf{x}
  =
  q(t)\int_\Omega g(\mathbf{x})v\,d\mathbf{x}.
\end{equation}

Используется пространство конечных элементов Лагранжа первой степени
на тетраэдральной сетке:
\[
  V_h = \{v_h:\; v_h|_K \in P_1(K)\}.
\]


\subsubsection{Полудискретная система}

После FEM-дискретизации получаем систему ОДУ:
\begin{equation}\label{eq:semi_discrete}
  \mathbf{M}\ddot{\mathbf{p}} + \mathbf{K}\mathbf{p}
  = q(t)\mathbf{F},
\end{equation}
где
\begin{equation}
  M_{ij} = \int_\Omega
  \frac{1}{\rho c^2}\,\phi_i\phi_j\,d\mathbf{x},
\end{equation}
\begin{equation}
  K_{ij} = \int_\Omega
  \frac{1}{\rho}\,\nabla\phi_j\cdot\nabla\phi_i\,d\mathbf{x},
\end{equation}
\begin{equation}
  F_i = \int_\Omega g(\mathbf{x})\phi_i\,d\mathbf{x}.
\end{equation}

Матрица жёсткости $\mathbf{K}$ собирается обычным образом через
\texttt{assemble\_matrix}. Для матрицы масс используется диагональное
приближение: вместо полной матрицы хранится вектор сумм по строкам.
Так явный шаг становится дешевле: не нужно решать линейную систему на
каждом временном слое.

\subsubsection{Явная схема по времени}

Для второй производной по времени используется центральная разность:
\begin{equation}
  \ddot{\mathbf{p}}^n \approx
  \frac{\mathbf{p}^{n+1}-2\mathbf{p}^{n}+\mathbf{p}^{n-1}}{\Delta t^2}.
\end{equation}
Подставляя это в~\eqref{eq:semi_discrete}, получаем явную схему:
\begin{equation}\label{eq:explicit_acoustic}
  \mathbf{p}^{n+1}
  =
  2\mathbf{p}^{n}-\mathbf{p}^{n-1}
  +
  \Delta t^2\,\mathbf{M}_L^{-1}
  \left(
    q(t_n)\mathbf{F}-\mathbf{K}\mathbf{p}^{n}
  \right).
\end{equation}

Начальные условия нулевые:
\[
  p(\mathbf{x},0)=0, \qquad
  \frac{\partial p}{\partial t}(\mathbf{x},0)=0.
\]

\subsubsection{Условие устойчивости}

Явная схема для волнового уравнения требует шага, согласованного
со скоростью распространения волны и минимальным размером элемента.
В коде шаг оценивается автоматически:
\begin{equation}
  \Delta t =
  \mathrm{safety}\,
  \frac{h_{\min}}{c_{\max}},
  \qquad
  \mathrm{safety}=0{,}25.
\end{equation}
Здесь $h_{\min}$ вычисляется через \texttt{cpp.mesh.h}.
\[
  c_{\max}=3{,}5\cdot 10^6~\text{мм/с}.
\]

\subsection{Реализация}

\subsubsection{Структура проекта}

Основные файлы:
\begin{itemize}
  \item \texttt{mini\_project/src/acoustic\_impulse.py}~--- основной
        акустический расчёт;
  \item \texttt{mini\_project/src/tools/preprocess.py}~--- загрузка VTU,
        перенос \texttt{Material ID}, создание DG0-полей;
  \item \texttt{mini\_project/src/tools/convert\_to\_vtu\_with\_materials.py}~---
        добавление физических параметров в VTU-файл;
\end{itemize}

Результаты пишутся в папку
\texttt{mini\_project/src/acoustic\_vtu\_frames\_other\_source/}
как последовательность \texttt{.vtu}-кадров.

\subsubsection{Контроль устойчивости}

Если шаг времени или амплитуда источника выбраны неправильно, решение
быстро начинает неограниченно расти, и расчёт останавливается.

Дополнительно задан порог
\[
  \max |p| \le 10^6 \text{Па}.
\]
Если давление превышает этот предел, программа сообщает, на каком шаге
произошла ошибка и останавливается.

\subsection{Результаты}

При запуске акустического расчёта импульс рождается около верхней границы
головы и распространяется через неоднородную анатомическую структуру.
Из-за различий в $c(\mathbf{x})$ и $\rho(\mathbf{x})$ фронт не остаётся
сферическим: на границах тканей возникают отражения и изменения длины
волны. Особенно заметное влияние дают костные области, где скорость
звука больше, чем в мягких тканях.

\begin{figure}[H]
  \centering
  \includegraphics[width=0.78\linewidth]{micro2/acoustic_impulse.png}
  \caption{Кадр визуализации акустического давления в ParaView.
           Красно-синие области соответствуют положительным и отрицательным
           фазам ультразвукового импульса.}
\end{figure}

По результатам расчёта были получены VTU-кадры и видеовизуализации.
См. \href{https://disk.360.yandex.ru/d/l9_8UlurSN1kMA}{примеры видео на Яндекс Диске}.

\subsection{Выводы}

В работе собран рабочий FEM-пайплайн для моделирования ультразвукового
импульса в 3D-модели головы. Главный результат~--- переход от готовой
анатомической VTU-сетки к физической DOLFINx-модели: метки материалов
сохранены, коэффициенты $\rho$ и $c$ назначены, а давление
рассчитано явной схемой на тетраэдральной сетке.

Физическая постановка соответствует линейной акустике в неоднородной
среде. Использование импульсного источника даёт короткий
ультразвуковой сигнал.
